\documentclass[12pt]{article}
\makeatletter
\def\input@path{{../../}}
\makeatother
\usepackage{sbc-template}
\makeatletter
\renewcommand{\inst}[1]{}
\makeatother
\usepackage{graphicx,url}
\usepackage[utf8]{inputenc}
\usepackage[T1]{fontenc}
\usepackage[brazil]{babel}
\usepackage{longtable,booktabs,array}
\providecommand{\tightlist}{%%
  \setlength{\itemsep}{0pt}\setlength{\parskip}{0pt}}
\sloppy
\title{Primeiro pipeline reproduzivel: previsao de demanda no Favorita com baselines e RC}
\author{}
\address{}
\begin{document}

\maketitle

\begin{resumo}
Depois de motivar o problema de negocio e introduzir Reservoir Computing, este artigo entrega o primeiro pipeline ponta a ponta da serie. Partimos do dataset Favorita, fixamos um recorte reproduzivel, construimos um baseline sazonal simples e comparamos esse baseline com um modelo de RC implementado no repositorio. O resultado nao e vendido como prova de superioridade do reservatorio. Pelo contrario: ele mostra por que um experimento inicial deve ser didatico, honesto e comparavel. No recorte \texttt{store\_nbr\ =\ 1}, \texttt{family\ =\ BEVERAGES}, com os ultimos 90 dias como teste, o Seasonal Naive obteve \texttt{MAE\ =\ 259.067}, \texttt{RMSE\ =\ 352.280}, \texttt{WAPE\ =\ 0.1196} e \texttt{sMAPE\ =\ 0.1361}, enquanto o RC obteve \texttt{MAE\ =\ 324.294}, \texttt{RMSE\ =\ 423.417}, \texttt{WAPE\ =\ 0.1498} e \texttt{sMAPE\ =\ 0.1653}. Em vez de frustrar o argumento da serie, esse resultado o fortalece: antes de pedir mais complexidade, o leitor aprende a construir um benchmark temporal limpo e a interpretar resultados sem autoengano.
\end{resumo}

\section{1. Introducao}\label{1-introducao}

Um erro comum em materiais introdutorios e apresentar o primeiro experimento apenas quando o metodo favorito do autor vence. Isso pode produzir um texto mais "bonito", mas produz um aprendizado pior. Neste artigo faremos o oposto: construir um pipeline minimo, inteiramente reproduzivel, e olhar o que ele realmente entrega.

Nosso objetivo aqui nao e provar que RC domina o problema de previsao de demanda. Nosso objetivo e ensinar:

\begin{itemize}
\tightlist
\item
  como escolher um recorte inicial do Favorita;
\item
  como organizar treino e teste sem vazamento;
\item
  como definir um baseline que merece respeito;
\item
  como executar um primeiro modelo de RC sobre o mesmo protocolo.
\end{itemize}

Com isso, o leitor ganha uma base metodologica solida para todo o restante da serie.

\section{2. Definicao do recorte inicial}\label{2-definicao-do-recorte-inicial}

O recorte experimental adotado nesta serie e:

\begin{itemize}
\tightlist
\item
  \texttt{store\_nbr\ =\ 1}
\item
  \texttt{family\ =\ BEVERAGES}
\item
  frequencia diaria
\item
  ultimos \texttt{90} dias reservados para teste
\end{itemize}

Essa escolha foi mantida em todas as implementacoes do repositorio para que os modelos sejam comparaveis. Trabalhar com uma unica serie nao esgota o Favorita, mas traz tres beneficios fundamentais:

\begin{enumerate}
\def\labelenumi{\arabic{enumi}.}
\tightlist
\item
  reduz custo computacional;
\item
  facilita interpretacao visual;
\item
  torna natural a comparacao entre familias de modelos muito diferentes.
\end{enumerate}

O primeiro pipeline da serie deve ser pequeno o suficiente para ser entendido por completo.

\section{3. Preparacao dos dados}\label{3-preparacao-dos-dados}

A preparacao adotada no projeto segue uma logica minimalista:

\begin{enumerate}
\def\labelenumi{\arabic{enumi}.}
\tightlist
\item
  filtrar loja e familia;
\item
  ordenar por data;
\item
  renomear o alvo como \texttt{y};
\item
  separar treino e teste por corte temporal;
\item
  construir atributos exogenos e ciclicos apenas a partir de informacoes disponiveis em cada instante.
\end{enumerate}

O ponto mais importante nao e o volume de preprocessing, mas a disciplina temporal. Em series temporais, um experimento simples e correto vale mais do que um pipeline sofisticado com vazamento.

No projeto, essa disciplina esta encapsulada em utilitarios compartilhados em \texttt{code/common/}.

\section{4. Baseline inicial: Seasonal Naive}\label{4-baseline-inicial-seasonal-naive}

O primeiro baseline escolhido e o Seasonal Naive com periodicidade semanal. Ele repete a estrutura recente de sete dias e serve como referencia minima forte para varejo diario.

Em termos simples, a previsao do dia \texttt{t} reaproveita o valor observado no dia \texttt{t\ -\ 7}, respeitando a intuicao de que segunda-feira se parece mais com segunda-feira do que com sabado.

Didaticamente, esse baseline e excelente porque:

\begin{itemize}
\tightlist
\item
  e facil de explicar;
\item
  e barato de executar;
\item
  captura um componente importante de sazonalidade;
\item
  obriga modelos mais sofisticados a mostrar que realmente estao agregando valor.
\end{itemize}

Os resultados obtidos no recorte da serie foram:

\begin{longtable}[]{@{}lrrrr@{}}
\toprule\noalign{}
Modelo & MAE & RMSE & WAPE & sMAPE \\
\midrule\noalign{}
\endhead
\bottomrule\noalign{}
\endlastfoot
Seasonal Naive & 259.067 & 352.280 & 0.1196 & 0.1361 \\
\end{longtable}

Esse e o primeiro numero que qualquer novo modelo precisa respeitar.

\section{5. Primeiro modelo de RC}\label{5-primeiro-modelo-de-rc}

O modelo de RC usado neste primeiro pipeline e uma versao de Echo State Network com:

\begin{itemize}
\tightlist
\item
  \texttt{n\_reservoir\ =\ 80}
\item
  \texttt{spectral\_radius\ =\ 0.6}
\item
  \texttt{leak\_rate\ =\ 0.15}
\item
  \texttt{washout\ =\ 14}
\item
  readout Ridge com \texttt{ridge\_alpha\ =\ 1e-3}
\end{itemize}

O vetor de entrada sequencial combina:

\begin{itemize}
\tightlist
\item
  demanda previa normalizada;
\item
  promocao normalizada;
\item
  componentes ciclicos do dia da semana;
\item
  componentes ciclicos do mes;
\item
  indicador de fim de semana.
\end{itemize}

O reservatorio atualiza um estado interno ao longo do tempo e o readout aprende a prever a demanda a partir desse estado expandido. A implementacao esta em \texttt{code/rc/model.py}.

\section{6. Resultados iniciais}\label{6-resultados-iniciais}

No mesmo protocolo experimental, o RC obteve:

\begin{longtable}[]{@{}lrrrr@{}}
\toprule\noalign{}
Modelo & MAE & RMSE & WAPE & sMAPE \\
\midrule\noalign{}
\endhead
\bottomrule\noalign{}
\endlastfoot
RC & 324.294 & 423.417 & 0.1498 & 0.1653 \\
\end{longtable}

Comparando baseline e RC:

\begin{longtable}[]{@{}lrrrr@{}}
\toprule\noalign{}
Modelo & MAE & RMSE & WAPE & sMAPE \\
\midrule\noalign{}
\endhead
\bottomrule\noalign{}
\endlastfoot
Seasonal Naive & 259.067 & 352.280 & 0.1196 & 0.1361 \\
RC & 324.294 & 423.417 & 0.1498 & 0.1653 \\
\end{longtable}

Neste primeiro recorte, o RC nao venceu o baseline. Essa constatacao e importante por duas razoes.

Primeiro, porque impede que o leitor associe "reservatorio" a uma promessa vazia de superioridade automatica. Segundo, porque abre perguntas muito mais interessantes do que uma simples tabela de vencedores:

\begin{itemize}
\tightlist
\item
  o reservatorio tem memoria adequada para esta serie?
\item
  o vetor de entrada esta informativo o bastante?
\item
  o baseline ja captura grande parte da sazonalidade relevante?
\item
  o ganho de RC apareceria com outro recorte, outra escala ou outro tuning?
\end{itemize}

Essas sao exatamente as perguntas que justificam o restante da serie.

\section{7. O que este primeiro experimento ensina}\label{7-o-que-este-primeiro-experimento-ensina}

Mesmo sem vencer, o primeiro pipeline entrega quatro aprendizados fundamentais.

\subsection{7.1 Baseline forte nao e detalhe}\label{71-baseline-forte-nao-e-detalhe}

Se o baseline for fraco, qualquer modelo sofisticado parece bom. Quando o baseline e serio, o experimento fica pedagogicamente util.

\subsection{7.2 RC nao deve ser tratado como magia}\label{72-rc-nao-deve-ser-tratado-como-magia}

Reservoir Computing e um paradigma interessante, mas ele nao elimina a necessidade de engenharia experimental. A qualidade do reservatorio, a representacao de entrada e o protocolo importam.

\subsection{7.3 Reproducibilidade e uma conquista em si}\label{73-reproducibilidade-e-uma-conquista-em-si}

Ao final deste artigo, o leitor ja possui um pipeline temporal reproduzivel, com recorte fixo, baseline, modelo alternativo e metricas comuns. Isso e muito mais valioso do que uma colecao de resultados desconectados.

\subsection{7.4 A derrota inicial organiza as proximas perguntas}\label{74-a-derrota-inicial-organiza-as-proximas-perguntas}

Os proximos artigos nao existem para "defender" RC. Eles existem para entender:

\begin{itemize}
\tightlist
\item
  que propriedades do reservatorio importam;
\item
  como comparar RC com outras familias de modelos;
\item
  em que sentido QRC herda o mesmo problema metodologico.
\end{itemize}

\section{8. Como este artigo se conecta ao benchmark maior}\label{8-como-este-artigo-se-conecta-ao-benchmark-maior}

O benchmark consolidado do projeto, apresentado mais adiante na serie, mostra que o RC tambem fica abaixo de ETS, XGBoost e Prophet neste recorte. Ainda assim, o modelo supera a LSTM simples implementada no projeto.

Isso e uma boa noticia didatica. O RC nao aparece como vencedor universal, nem como fracasso irrelevante. Ele aparece como uma alternativa temporal com estrutura propria, que merece ser compreendida e comparada.

\section{9. Conclusao}\label{9-conclusao}

O primeiro pipeline reproduzivel da serie cumpre sua funcao. Ele mostra que:

\begin{itemize}
\tightlist
\item
  o problema pode ser colocado em pe de forma limpa;
\item
  um baseline simples pode ser muito competitivo;
\item
  o RC ja pode ser executado de ponta a ponta no Favorita;
\item
  desempenho menor que o baseline nao invalida o metodo, mas qualifica a discussao.
\end{itemize}

No proximo artigo, vamos investigar o mecanismo por tras do RC: memoria, nao linearidade e capacidade de representacao.

\section{Entregaveis associados no repositorio}\label{entregaveis-associados-no-repositorio}

\begin{itemize}
\tightlist
\item
  baseline sazonal: \texttt{code/baselines/seasonal\_naives/}
\item
  RC classico: \texttt{code/rc/}
\item
  benchmark consolidado: \texttt{code/RESULTS.md}
\end{itemize}

\section{Referencias}\label{referencias}

\begin{enumerate}
\def\labelenumi{\arabic{enumi}.}
\tightlist
\item
  Jaeger, H. The "echo state" approach to analysing and training recurrent neural networks.
\item
  Hyndman, R. J., Athanasopoulos, G. Forecasting: Principles and Practice.
\item
  Kaggle. Corporacion Favorita Grocery Sales Forecasting.
\end{enumerate}

\section*{Resultados Computacionais Associados}

Os resultados computacionais associados a este artigo estao organizados na subpasta \texttt{computational\_results\_20260402\_222902/}, localizada dentro desta pasta \LaTeX. Esse diretorio contem o arquivo \texttt{REPORT.md}, tabelas em CSV e imagens comparativas apropriadas ao escopo do artigo.

\end{document}
